\section{Experiments}
\label{sec:experiments}

We perform a comprehensive empirical evaluation of HessMerge against established post-hoc model merging baselines. We evaluate joint multi-task performance under six diverse downstream deployment quantization schemas using a Vision Transformer backbone across four visual classification domains.

\subsection{Experimental Setup}
\paragraph{Backbone Architecture and Expert Training.} 
We adopt the Vision Transformer (\texttt{vit\_tiny\_patch16\_224}) \cite{vit} from the \texttt{timm} library \cite{timm} as our shared pretrained base model $\mathcal{M}_{\text{pre}}$. This lightweight backbone contains approximately 5.7M parameters grouped into $L=14$ layer-wise modules: the embedding layer (including patch projection, class token, and positional embeddings) as Layer 0, twelve consecutive Transformer blocks as Layers 1--12, and the final layer normalization layer as Layer 13.

We fine-tune the shared pretrained model independently on four distinct datasets: \textbf{MNIST} \cite{mnist}, \textbf{FashionMNIST} \cite{fashionmnist}, \textbf{CIFAR-10} \cite{cifar10}, and \textbf{SVHN} \cite{svhn}. Each expert is fine-tuned for exactly 3 epochs using the AdamW optimizer with a learning rate of $10^{-3}$, a weight decay of $0.01$, and a batch size of $64$, without data augmentation. For SVHN, this rapid fine-tuning represents a challenging, noisy expert case (achieving $\approx 18.4\%$ accuracy due to dataset complexity), which serves as an excellent stress test for merging robustness. The converged individual expert test-set accuracies (representing task-specific performance ceilings) are: MNIST (\textbf{23.3\%}), FashionMNIST (\textbf{41.7\%}), CIFAR-10 (\textbf{29.6\%}), and SVHN (\textbf{18.4\%}). The task classification heads are kept task-specific and swapped dynamically during evaluation, while the shared backbone parameters are merged.

\paragraph{Calibration Stream and Test-Time Adaptation.}
For test-time adaptation, we construct an extremely lightweight calibration stream containing exactly $B=16$ unlabeled test-set queries per task, yielding a joint calibration dataset of size $N=64$. Optimization is conducted for $T=100$ steps using the Adam optimizer with a learning rate of $\eta = 10^{-2}$. For HessMerge, we set the Hessian curvature penalty coefficient to $\gamma = 0.5$.

\subsection{Comparative Baselines}
We evaluate HessMerge against five prominent weight-space model merging frameworks:
\begin{enumerate}
    \item \textbf{Uniform Task Arithmetic (No TTA):} A static, zero-compute baseline blending model weights with uniform coefficients $\lambda_k^l = 1/K = 0.25$ across all layers.
    \item \textbf{AdaMerging (Unregularized TTA) \cite{adamerging}:} The standard adaptive merging scheme, optimizing layer-wise coefficients $\Lambda \in \mathbb{R}^{L \times K}$ purely on the unsupervised multi-task entropy loss without any regularization.
    \item \textbf{RegCalMerge \cite{zipmerge}:} An adaptive scheme incorporating Elastic Spatial Regularization (ESR), implemented as a Total Variation (TV) penalty over adjacent layers to smooth coefficient transitions: $\mathcal{R}_{\text{TV}}(\Lambda) = \beta \sum_k \sum_l (\lambda_k^{l+1} - \lambda_k^l)^2$.
    \item \textbf{Q-Merge:} A quantization-aware adaptive merging scheme simulating INT8 symmetric tensor-wise quantization in the forward pass of the optimizer and utilizing Straight-Through Estimators (STE) to propagate gradients.
    \item \textbf{PolyMerge:} A subspace-constrained adaptive scheme that parameterizes the layer-wise coefficients as a continuous low-degree polynomial of network depth: $\lambda_k^l = \sigma\left(a_k + b_k (l/L) + c_k (l/L)^2\right)$, optimizing only the polynomial coefficients.
\end{enumerate}

\subsection{Deployment Quantization Schemas}
To thoroughly evaluate robustness under realistic edge hardware deployment, we pass the merged continuous parameters through six distinct Post-Training Quantization (PTQ) target schemas before evaluation:
\begin{itemize}
    \item \textbf{FP32 (No Quantization):} The unquantized, baseline floating-point model.
    \item \textbf{INT8 Uniform Symmetric (Tensor-wise):} Symmetric quantization mapped to $[-127, 127]$, with a single scale computed over the entire weight tensor.
    \item \textbf{INT8 Uniform Symmetric (Channel-wise):} Symmetric quantization with independent scales computed for each output channel.
    \item \textbf{INT8 Uniform Asymmetric (Tensor-wise):} Asymmetric quantization mapped to $[-128, 127]$, with a scale and zero-point computed over the entire weight tensor.
    \item \textbf{INT8 Uniform Asymmetric (Channel-wise):} Asymmetric quantization with independent scales and zero-points computed for each output channel.
    \item \textbf{INT4 Uniform Symmetric (Channel-wise):} Aggressive low-precision symmetric quantization mapped to $[-7, 7]$ with independent per-channel scales.
\end{itemize}
All task classification heads are kept in FP32 format during evaluation to isolate backbone representation robustness.
